\scanpage{7}
\uline{connexe} \uline{orienté} ``revient au $=$'' que de se donner un $n\in\mathbb{N}$, $n\geq2$, et un \inserted{tel} torseur sous $\mathbb{Z}/n\mathbb{Z}$ (savoir $S$, \unsure{ou|--} $A$ \unsure{au|--} choix -- NB $\unsure{r|\sigma}:S\longrightarrow A$ est compatible avec l'action de $G^{\unsure{0|\circ}}=\mathbb{Z}/n\mathbb{Z}$ i.e. avec l'action de $u_{\omega}$.

Se donner un \uline{contour connexe} équivaut : la donnée de

$1^{\circ})$ un \unsure{gpe|grpe} cyclique $G^{0}$ (d'ordre $n\geq2$)

$2^{\circ})$ une \deleted{\unsure{ens|un}} partie de $G^{0}$ \deleted{formée d'} de \unsure{la|2} forme $\{g,\unsure{s|g}^{-1}\}$, où $g$ est un générateur \inserted{(\unsure{donc|des} partie de card $2$ \unsure{sss|ss|ssi} $n\ne2$)}

$3^{\circ})$ Un torseur $S$ sous $G^{0}$

$\bullet$ Du moins si $n\geq3$ \deleted{\unsure{sinon|si on}} ; pour une description

$\bullet$ générale, il faut de plus se donner, en lieu du $2^{\circ}$, un

$2$ bis) un ens $\Omega$ à deux \unsure{éléts|élémts|éléments}, et une application $\Omega\xrightarrow{i}G$ telle que, si $\Omega=\{a,b\}$, on ait $i(a)i(b)=1$, $i(a)$ gén. de $G$

% Short rising separator stroke.
\figure{\begin{tikzpicture}[x=1cm,y=1cm]
\draw[line width=.4pt] (0,0) -- (1,.11);
\end{tikzpicture}}

Contour orienté (pas néc. \unsure{connexe|connexe.}) $=$ ensemble sur lequel $\mathbb{Z}$ opère (i.e \unsure{avec|ou} autom \unsure{$u$|$n$})

Contour (pas néc connexe) $=\unsure{\ldots|un}$
% archive sheet number: 6
